% ======================================================================
%  New short subsection / paragraph: what each Hermite scale contributes.
%
%  Place wherever the Hermite bank and its four scales are introduced
%  (Chapter 3, or Section 4.4 if the bank is restated there). It is a
%  measured property of the feature, not a result of the ablation, so it
%  does not belong in 4.7.
%
%  SOURCE OF THE NUMBERS: the per-channel residual RMS printed at the first
%  iteration of the V3 nominal run,
%      chan 0 = 1115.35, chan 1 = 1445.22, chan 2 = 1315.24, chan 3 = 464.302
%  squared and normalized to sum to unity. The feature is channel-blocked,
%  264000 = 4 x 66000 components.
%
%  %%% VERIFY: the association of channel index with scale assumes the
%  %%% feature vector stacks the responses in the order the images are
%  %%% passed to buildFrom, i.e. (w_a, w_h, w_v, w_d) = (s1, s2, s3, s4).
%  %%% Confirm against the loop in the customized vpFeatureLuminance before
%  %%% publishing the sigma column. The shares themselves do not depend on
%  %%% this; only the labelling of the rows does.
% ======================================================================

\paragraph{Contribution of the individual scales.}
The four scales of the bank do not contribute equally to the feature error.
Measuring the residual separately in each of the four blocks of the stacked
feature vector at the start of a nominal run gives the distribution of
Table~\ref{tab:scale-energy}. Two points follow. First, the dependence on
$\sigma$ is not monotone: the response is largest at $\sigma_2=1.1$ rather
than at the finest scale, because a $9\times9$ support samples a
$\sigma=0.5$ kernel too coarsely for its taps to contribute coherently, so
the finest channel behaves closer to an identity than to a filter. Second,
the coarsest scale is nearly redundant, carrying $4\%$ of the residual
energy against $33$--$40\%$ for the two intermediate scales; at
$\sigma_4=2.8$ the smoothing has removed most of the spatial detail that
carries pose information.

This is the quantitative justification for treating the scales
$\{\sigma_k\}$ as parameters to be chosen rather than as a fixed geometric
progression: a bank whose members contribute between $4\%$ and $40\%$ of the
usable signal is demonstrably unbalanced, and the choice of scales is
therefore a legitimate target for optimization rather than a detail of
implementation.

\begin{table}[t]
    \centering
    \caption{Share of the total feature-error energy carried by each scale of
    the Hermite bank, measured at the initial pose of the nominal run. The
    feature has $264\,000$ components, four blocks of $66\,000$ over the
    $220\times300$ interior region left by the filter border.}
    \label{tab:scale-energy}
    \small
    \begin{tabular}{@{}clrr@{}}
        \toprule
        Channel & Scale & Residual RMS & Share of energy \\
        \midrule
        0 & $\sigma_1=0.5$ & $1115.4$ & $23.6\%$ \\
        1 & $\sigma_2=1.1$ & $1445.2$ & $39.6\%$ \\
        2 & $\sigma_3=1.8$ & $1315.2$ & $32.8\%$ \\
        3 & $\sigma_4=2.8$ & $\phantom{0}464.3$ & $\phantom{0}4.1\%$ \\
        \midrule
        \multicolumn{2}{@{}l}{Whole feature} & $1148.7$ & $100\%$ \\
        \bottomrule
    \end{tabular}
\end{table}
